\section{Fibrations with small fibre Chow groups: a geometric reduction
of the cycle
gate}\label{fibrations-with-small-fibre-chow-groups-a-geometric-reduction-of-the-cycle-gate}

\textbf{Scope.} This is a deduction from Vial's published motivic
fibration theorem, with rational coefficients. It constructs a Chow
decomposition of the original variety. In the cases below where the
remaining summand comes from a curve or a finite étale scheme, the
forward-cycle existence gate in CC-XC15 is therefore discharged. No
claim about an arbitrary fibration is made from its generic fibre alone.

Write \(\mathcal C_{\mathrm{ab,Art}}(k)\) for the full rational Chow
subcategory generated by abelian varieties, finite étale schemes and
Tate twists, closed under sums, tensor products, duals and summands.
Smooth projective curves belong to this category: over \(\bar k\) this
is the curve/abelian equivalence in Vial (2017, p.~195), and the actual
restriction--corestriction proof in CC-XC16 descends this membership to
\(k\). Write \(\mathcal A_{q,E}\) for the fully embedded abelian--Artin
subcategory of the \(E=\overline{\mathbf Q}_\ell\)-linear homological
category used in CC-XC01--15. Thus certification means membership in
\(\mathcal A_{q,E}\), with its already-proved rational, all-\(\ell\)
consequences.

\subsection{\texorpdfstring{CC-XC18 --- A residual motive of dimension
\(\dim X-2n\)}{CC-XC18 --- A residual motive of dimension \textbackslash dim X-2n}}\label{cc-xc18-a-residual-motive-of-dimension-dim-x-2n}

Let \(k=\mathbf F_q\). Let \(f:X\to B\) be a flat, projective,
surjective morphism of smooth projective geometrically integral
\(k\)-varieties. Put

\[
       D=\dim X,\qquad b=\dim B\ge1,\qquad r=D-b\ge1.
 \tag{CC-XC.Q1}
\]

Let \(\Omega\supset\bar k\) be a universal domain. Suppose \(n\) is an
integer with \(1\le 2n\le D\) and

\[
 \dim_{\mathbf Q}\mathrm{CH}_j(X_t)_{\mathbf Q}=1
 \quad\text{for every }t\in B(\Omega)\text{ and every }0\le j<n.
 \tag{CC-XC.Q2}
\]

The notation includes singular and nonreduced fibres; Chow groups use
their usual scheme-theoretic cycles modulo rational equivalence. Then
there are a smooth projective \(k\)-variety \(Z\) of dimension \(D-2n\),
an idempotent \(e\in\mathrm{CH}^{\dim Z}(Z\times Z)_{\mathbf Q}\), and
an isomorphism of rational Chow motives

\[
 h(X)\simeq
 \left(\bigoplus_{i=0}^{r} h(B)(-i)\right)
       \oplus (Z,e)(-n).
 \tag{CC-XC.Q3}
\]

Here and below twists are in the contravariant convention of this
project. Consequently:

\begin{enumerate}
\def\labelenumi{\arabic{enumi}.}
\tightlist
\item
  \(h(X)\in\mathcal C_{\mathrm{ab,Art}}(k)\) if and only if both
  \(h(B)\) and \((Z,e)\) belong to that category.
\item
  \(X\) is certified if and only if \(B\) is certified and the
  homological image of \((Z,e)\) belongs to \(\mathcal A_{q,E}\).
\item
  If \(D-2n\le1\), the extra condition on \((Z,e)\) is automatic. Thus
  in this case Chow membership of \(X\) is equivalent to Chow membership
  of \(B\), and certification of \(X\) is equivalent to certification of
  \(B\).
\end{enumerate}

\textbf{Complete deduction.} Formula (CC-XC.Q3) is exactly Vial (2013,
Theorem 4.2, p.~1534), translated from his covariant convention. His
positive Lefschetz twist becomes our negative Tate twist; the idempotent
is transposed under the corresponding opposite-category convention, and
we have denoted the resulting idempotent by \(e\). The cited theorem
assumes (CC-XC.Q2) for every \(\Omega\)-point, and does not impose
characteristic zero, smoothness of \(f\), or a smooth discriminant. The
projective-surjective convention is fixed at the beginning of the paper,
p.~1522. We impose \(2n\le D\) to keep the displayed residual dimension
nonnegative.

The isomorphism is at the level of Chow correspondences. Its proof on
p.~1534 uses the graph/hyperplane splitting of Theorem 1.4 and the
self-dual base projector of Proposition 1.5. The fibre-Chow hypothesis
makes both low-dimensional Chow groups of the residual projector and of
its transpose vanish, and Proposition 4.1 then gives the stated residual
variety. We invoke that established theorem; equality of zeta functions
is not being used to infer (CC-XC.Q3).

Both categories in assertions 1 and 2 are closed under finite sums,
summands and Tate twists. If the two indicated objects belong,
(CC-XC.Q3) therefore puts \(h(X)\) in the category. Conversely, \(h(B)\)
is the \(i=0\) direct summand of (CC-XC.Q3), and \((Z,e)\) is a twist of
another direct summand, so membership of \(h(X)\) forces both. Passing
the actual Chow isomorphism to homological motives is legitimate without
assuming numerical--homological comparison for \(Z\): the latter
category is defined using the kernel of the chosen cycle realization.

If \(\dim Z=0\), smoothness and projectivity over a finite field imply
that \(Z\) is finite étale. If \(\dim Z=1\), its components are smooth
projective curves, possibly with larger finite constant fields; the
finite-field descent argument in CC-XC16 includes those components and
their Artin factors. Thus \(h(Z)\), and hence its summand \((Z,e)\),
lies in \(\mathcal C_{\mathrm{ab,Art}}(k)\) whenever \(\dim Z\le1\).
This proves assertion 3.

When certification follows, CC-XC14 constructs the fixed rational target
and CC-XC15's forward-only search terminates. This is a termination
consequence of actual cycle existence, with no claimed degree or runtime
bound. The theorem does not assert an effective construction of the
auxiliary \(Z\) from the fibre-Chow hypothesis alone. \(\square\)

\subsection{Quadric-fibration consequences of
CC-XC18}\label{quadric-fibration-consequences-of-cc-xc18}

Suppose every geometric fibre of \(f\) is a quadric hypersurface of
dimension \(r\). Retain flatness, smooth projective \(X,B\), and
projectivity and surjectivity of \(f\); the fibres themselves need not
be smooth. The quadric case of Vial (2013, Corollary 4.4 and its proof,
p.~1535) supplies (CC-XC.Q2) with

\[
 n=\left\lceil\frac r2\right\rceil,\qquad
 \dim Z=
 \begin{cases}
 b,&r\text{ even},\\
 b-1,&r\text{ odd}.
 \end{cases}
 \tag{CC-XC.Q4}
\]

This primary result is stated over an arbitrary field, with rational
Chow groups. It therefore applies in characteristic two as well as in
odd characteristic. We use its quadric case and do not impose additional
discriminant or corank hypotheses.

The following are exact consequences.

\begin{longtable}[]{@{}
  >{\raggedright\arraybackslash}p{(\columnwidth - 4\tabcolsep) * \real{0.3333}}
  >{\raggedright\arraybackslash}p{(\columnwidth - 4\tabcolsep) * \real{0.3333}}
  >{\raggedright\arraybackslash}p{(\columnwidth - 4\tabcolsep) * \real{0.3333}}@{}}
\toprule\noalign{}
\begin{minipage}[b]{\linewidth}\raggedright
Base and relative dimension
\end{minipage} & \begin{minipage}[b]{\linewidth}\raggedright
Remaining Chow summand
\end{minipage} & \begin{minipage}[b]{\linewidth}\raggedright
Consequence for the original \(X\)
\end{minipage} \\
\midrule\noalign{}
\endhead
\bottomrule\noalign{}
\endlastfoot
\(b=1\), any \(r\ge1\) & A summand of a curve motive if \(r\) is even,
or of a finite étale motive if \(r\) is odd &
\(h(X)\in\mathcal C_{\mathrm{ab,Art}}(k)\); its forward-cycle search
terminates. \\
\(b=2\), \(r\) odd & A twisted summand of a curve motive & Chow
membership of \(X\) is equivalent to that of \(B\); certification of
\(X\) is equivalent to that of \(B\). \\
\(b=2\), \(r\) even & One specified twisted projector summand of a
surface motive & The exact remaining gate is membership of \(h(B)\) and
of \((Z,e)\), not necessarily of the whole surface \(Z\). \\
General \(b\ge1\) & Dimension at most \(b\), and \(b-1\) for odd \(r\) &
The high-dimensional forward-cycle problem reduces to the base and this
single residual projector. \\
\end{longtable}

\textbf{Verification of these consequences.} For an integer \(j\),
\(j<r/2\) is equivalent to \(0\le j<\lceil r/2\rceil\) once the lower
bound \(j\ge0\) is included. Substituting this \(n\) into
\(D-2n=b+r-2\lceil r/2\rceil\) gives (CC-XC.Q4). Because \(b\ge1\), the
residual dimension is nonnegative and the range imposed in CC-XC18 is
satisfied. The first two rows now follow from assertion 3 and the
membership of curves. The third and fourth rows are precisely assertions
1 and 2, retaining the projector \(e\) instead of replacing it by an
unwarranted statement about all of \(h(Z)\).

In the first row, \(h(B)\) and the residual motive are generated by
curves and Artin motives. Their twists in (CC-XC.Q3) therefore prove
actual Chow inclusion, not just Kimura finite-dimensionality. In the
second row, if \(B\) is only known to be certified, the conclusion is
certification of \(X\); a Chow-membership assertion for \(B\) is not
inserted. Examples with a Chow abelian--Artin base include an abelian
surface and the universal-zero-cycle surfaces of CC-XC17. \(\square\)

\subsection{Scope of the geometric
input}\label{scope-of-the-geometric-input}

The theorem uses all geometric fibres; we have not replaced that
condition by a condition on only the generic fibre. Vial explicitly
distinguishes these hypotheses in the introduction, pp.~1523--1524,
including the change of Chow groups in singular fibres of a conic
bundle. In the quadric case the required lower Chow groups remain
one-dimensional for the entire allowed fibre class, as encoded in
Corollary 4.4; this is why degenerations cause no extra hypothesis in
the conclusion above.

The residual surface in the even-dimensional quadric case is left
visible. A Chow--Künneth decomposition or Kimura finite-dimensionality
alone is not used as a substitute for abelian-sector membership. Nor do
the above statements claim that numerical and rational equivalence agree
on all products of \(X\). The conclusions needed for cote 15 are the
stated Chow inclusion or the exact lower-dimensional membership
reduction, followed by the existing finite forward-cycle search.

\subsection{Primary bibliography and
locators}\label{primary-bibliography-and-locators}

\begin{itemize}
\tightlist
\item
  Vial, C. (2013). Algebraic cycles and fibrations. \emph{Documenta
  Mathematica, 18}, 1521--1553. https://doi.org/10.4171/DM/435 .
  Publisher PDF:
  https://ems.press/content/serial-article-files/26216?nt=1 . Exact
  locators: standing hypotheses and rational coefficients, pp.~1522 and
  1525--1526; Theorem 1.4 and Proposition 1.5, pp.~1528--1529; Theorem
  4.2 and its proof, p.~1534; Corollary 4.4 and its proof, p.~1535. The
  publisher version was received 23 March 2012 and revised 19 July 2013.
  Some circulated files have obsolete pagination; the locators here
  refer to the publisher's 1521--1553 version.
\item
  Vial, C. (2017). Remarks on motives of abelian type. \emph{Tohoku
  Mathematical Journal, Second Series, 69}(2), 195--220.
  https://doi.org/10.2748/tmj/1498269623 . Exact input: p.~195,
  generation by curves and abelian varieties, together with the
  finite-field Chow descent proved in CC-XC16. Published source:
  https://www.jstage.jst.go.jp/article/tmj/69/2/69\_1498269623/\_pdf/-char/en
  .
\end{itemize}
